\documentclass[10pt]{article}
\usepackage[letterpaper,top=0.5in,bottom=0.5in,left=1in,right=1in]{geometry}
\usepackage{booktabs}
\usepackage{amsmath}
\usepackage{array}
\usepackage{siunitx}

\pagestyle{empty}

% Consistent number/unit formatting across the whole document
\sisetup{
	separate-uncertainty = true,
	table-format         = 2.2,
	detect-weight        = true,
	detect-family        = true,
}

% Setup table space aboth and below tables
\setlength{\abovetopsep}{16pt}
\setlength{\belowbottomsep}{16pt}

\newcommand{\fivexfive}{\ensuremath{5\times5}}
\newcommand{\mgdrop}{\ensuremath{\mathrm{mg/drop}}}

\begin{document}

\begin{center}
	{\LARGE\bfseries Fichot \fivexfive{} Quicksheet}
	\vspace{2pt}
	{\hrule height 0.8pt}
	\vspace{5pt}
\end{center}

This sheet is for usage with the Alain Fichot \fivexfive{} glaze testing fixture. Each well on
my fixture is measured at \qty{2.2}{cm} $\times$ \qty{2.4}{cm} and \qty{1}{cm} deep. The section
for Glaze Slurry contains information for the Leach 4321 glaze at ${\sim}\,\num{1.45}$ specific
gravity. The section on Additions contains information on how to create bottles
of addition material, generally a single colorant oxide, used to add ${\sim}X\%$ of
colorant per drop.

\section*{Glaze Slurry}

Dose each well of the Fichot \fivexfive{} fixture with a volume of glaze slurry from the
table below. The reference amount for most glazes is \qty{0.8}{mL} at ${\sim}\,\qty{0.4}{mm}$ thick,
which is what the standard addition amounts are based on below.

\begin{center}
	\begin{tabular}{
			S[table-format=1.1]
			S[table-format=1.2]
			S[table-format=1.2]
			S[table-format=1.2]
			S[table-format=1.2]
		}
		\toprule
		{Volume (\unit{mL})} & {Fired Thickness (\unit{mm})} & {Specific Gravity} & {Wet Mass (\unit{g})} & {Dry Mass (\unit{g})} \\
		\midrule
		0.6                  & 0.30                          & 1.44               & 0.87                  & 0.43                  \\
		0.8                  & 0.39                          & 1.44               & 1.16                  & 0.58                  \\
		1.0                  & 0.49                          & 1.44               & 1.44                  & 0.72                  \\
		1.2                  & 0.59                          & 1.44               & 1.73                  & 0.87                  \\
		\bottomrule
	\end{tabular}
\end{center}

\section*{Additions}

To make a standard addition bottle add \qty{3}{g} of colorant (representing ${\sim}\,\qty{0.5}{\percent}$ colorant
per \qty{0.05}{mL} drop) to \qty{50}{g} of distilled water. Optionally, add \qty{0.75}{g}
of sodium bentonite for better colorant suspension. Mix well and let sit at least 24 hours
before use for materials to fully hydrate.

To verify the colorant solution is accurate, weigh 100 drops of the mixture in grams ($W$),
let dry completely, and weigh the residue in grams ($R$). Use the formulas below to compute
\mgdrop, then convert to percentage per drop with the Dry Mass value above:

\vspace{2pt}
\begin{center}
	\setlength{\arraycolsep}{2pt}
	\renewcommand{\arraystretch}{2.2}
	\begin{tabular}{l l}
		\textbf{Colorant amount in $grams$:}  & $\displaystyle \mgdrop = 10\,R$                                       \\
		\textbf{...with $X grams$ bentonite:} & $\displaystyle \mgdrop = 10\,R \times \frac{m_c}{m_c + X}$            \\
		\textbf{Percentage Conversion:}       & $\displaystyle \%/drop = \frac{\mgdrop}{10 \times (\text{Dry Mass})}$ \\
	\end{tabular}

	\vspace{4pt}
	{\small$R$ = dry residue from 100 drops (\unit{g})\quad}

	{\small$m_c$ = colorant added to the bottle (\unit{g})\quad}
\end{center}

\vspace{4pt}

The bottle additions are targeted to add 0.5\% of colorant per drop based on the
\qty{0.4}{mm} thick reference glaze. There are cases where one might want to use
existing colorant bottles for other volumes of glaze. Following is a table that
shows percentages for other amounts of glaze based on the dry mass of the Glaze
Slurry. Most glazes will scale linearly with specific gravity (SG). Following is
the formula one can use to find dose volume for a given thickness with a new
specific gravity:

\begin{center}

	$\text{Dry Mass (g)} = 1.64 \times \text{Volume (mL)} \times (\text{SG} - 1)$

	\setlength{\tabcolsep}{11pt}
	\begin{tabular}{
			l
			S[table-format=1.2]
			S[table-format=1.2]
			S[table-format=1.2]
			S[table-format=1.2]
		}
		\toprule
		                                 & \multicolumn{4}{c}{Glaze Dry Mass (\unit{g})}                                                                      \\
		\cmidrule(lr){2-5}
		{Reference \%/drop}              & {0.43}                                        & {0.58}               & {0.72}               & {0.87}               \\
		\midrule
		\qty{0.2}{\percent} (\qty{1}{g}) & \qty{0.23}{\percent}                          & \qty{0.17}{\percent} & \qty{0.14}{\percent} & \qty{0.11}{\percent} \\
		\qty{0.5}{\percent} (\qty{3}{g}) & \qty{0.70}{\percent}                          & \qty{0.52}{\percent} & \qty{0.42}{\percent} & \qty{0.34}{\percent} \\
		\qty{1}{\percent} (\qty{6}{g})   & \qty{1.40}{\percent}                          & \qty{1.03}{\percent} & \qty{0.83}{\percent} & \qty{0.69}{\percent} \\
		\bottomrule
	\end{tabular}
\end{center}


\end{document}
